\documentclass[12pt]{article}
\usepackage[utf8]{inputenc}
\usepackage{amsmath, amssymb, amsthm}
\usepackage{geometry}
\usepackage{booktabs}
\usepackage{algorithm}
\usepackage{algpseudocode}
\geometry{margin=1in}

\title{Problem 41: Pandigital Prime}
\author{Project Euler Solutions}
\date{}

\newtheorem{theorem}{Theorem}
\newtheorem{lemma}[theorem]{Lemma}
\newtheorem{proposition}[theorem]{Proposition}
\newtheorem{definition}[theorem]{Definition}
\newtheorem*{remark}{Remark}

\begin{document}
\maketitle

\section{Problem Statement}

We say that an $n$-digit number is \emph{pandigital} if it uses each of the digits $1, 2, \ldots, n$ exactly once. Find the largest $n$-digit pandigital number that is also prime.

\section{Definitions}

\begin{definition}
For $n \in \{1, 2, \ldots, 9\}$, an \emph{$n$-pandigital number} is a positive integer whose decimal representation is a permutation of $\{1, 2, \ldots, n\}$.
\end{definition}

\begin{definition}
The \emph{digit sum} of the $n$-pandigital class is $S(n) = \sum_{k=1}^{n} k = \frac{n(n+1)}{2}$.
\end{definition}

\section{Mathematical Development}

\begin{theorem}[Divisibility by 3]
An $n$-pandigital number is divisible by $3$ if and only if $3 \mid S(n)$. In particular, every $n$-pandigital number is divisible by $3$ for $n \in \{2, 3, 5, 6, 8, 9\}$.
\end{theorem}

\begin{proof}
Let $N$ be an $n$-pandigital number with digits $a_1, \ldots, a_n$. By the divisibility rule for~3:
\[
N \equiv \sum_{i=1}^n a_i \pmod{3}.
\]
Since $\{a_1, \ldots, a_n\}$ is a permutation of $\{1, \ldots, n\}$, the digit sum equals $S(n) = n(n+1)/2$ for every $n$-pandigital number. Hence $3 \mid N$ if and only if $3 \mid S(n)$.

Computing $S(n) \bmod 3$:
\begin{center}
\begin{tabular}{c|ccccccccc}
\toprule
$n$ & 1 & 2 & 3 & 4 & 5 & 6 & 7 & 8 & 9 \\
\midrule
$S(n)$ & 1 & 3 & 6 & 10 & 15 & 21 & 28 & 36 & 45 \\
$S(n) \bmod 3$ & 1 & 0 & 0 & 1 & 0 & 0 & 1 & 0 & 0 \\
\bottomrule
\end{tabular}
\end{center}
For $n \in \{2,3,5,6,8,9\}$, we have $3 \mid S(n)$, so every $n$-pandigital number is divisible by~3. Since every such number exceeds 3, it is composite.
\end{proof}

\begin{lemma}
No $1$-pandigital prime exists.
\end{lemma}

\begin{proof}
The unique $1$-pandigital number is $1$, which is not prime.
\end{proof}

\begin{theorem}[Admissible values of $n$]
The only values of $n$ for which an $n$-pandigital prime can exist are $n = 4$ and $n = 7$.
\end{theorem}

\begin{proof}
By Lemma~2, $n = 1$ is excluded. By Theorem~1, $n \in \{2,3,5,6,8,9\}$ forces every $n$-pandigital number to be composite. The remaining candidates are $n = 4$ and $n = 7$, for which $S(n) \equiv 1 \pmod{3}$.
\end{proof}

\begin{remark}
Theorem~2 does not guarantee existence of pandigital primes for $n = 4$ or $n = 7$; it establishes these as the only admissible values. Existence is verified computationally in Theorem~3.
\end{remark}

\begin{theorem}
The largest pandigital prime is $7652413$.
\end{theorem}

\begin{proof}
By Theorem~2, it suffices to search among $7$- and $4$-pandigital numbers. Since every $7$-pandigital number exceeds every $4$-pandigital number ($1234567 > 4321$), we search $7$-pandigital numbers first in descending lexicographic order.

There are $7! = 5040$ permutations of $\{1,2,\ldots,7\}$. Enumerating in descending order and testing primality, the first prime encountered is $7652413$.

\emph{Verification.} We have $\lfloor \sqrt{7652413} \rfloor = 2766$. By trial division, $p \nmid 7652413$ for every prime $p \leq 2766$, confirming primality.

Since $7652413$ is the first prime in the descending enumeration of $7$-pandigital numbers, it is the largest $7$-pandigital prime, hence the largest pandigital prime overall.
\end{proof}

\section{Algorithm}

\begin{algorithm}[H]
\caption{Largest Pandigital Prime}
\begin{algorithmic}[1]
\For{$n \in \{7, 4\}$} \Comment{Theorem~2: only admissible values}
    \State $\text{digits} \gets \{1, 2, \ldots, n\}$
    \For{each permutation $P$ of digits in descending lexicographic order}
        \State $N \gets \textsc{DigitsToInteger}(P)$
        \If{$\textsc{IsPrime}(N)$}
            \State \Return $N$
        \EndIf
    \EndFor
\EndFor
\end{algorithmic}
\end{algorithm}

\section{Complexity Analysis}

\begin{proposition}[Time complexity]
The algorithm terminates in $O(n! \cdot \sqrt{N_{\max}})$ time in the worst case, where $n \leq 7$ and $N_{\max} = 7654321$.
\end{proposition}

\begin{proof}
In the worst case, all $7! = 5040$ permutations are enumerated. Each primality test costs $O(\sqrt{N_{\max}})$ by trial division, where $\sqrt{N_{\max}} < 2767$. Total:
\[
O(7! \cdot \sqrt{7654321}) = O(5040 \times 2767) \approx O(1.39 \times 10^7).
\]
\end{proof}

\begin{proposition}[Space complexity]
The algorithm uses $O(n)$ auxiliary space, where $n \leq 7$.
\end{proposition}

\begin{proof}
The only data structures are the current permutation and loop variables.
\end{proof}

\section{Answer}

\[
\boxed{7652413}
\]

\end{document}
